% ============================================
% Section: Hardware Design Module
% ============================================
\label{sec:hardware}

Co-LIC2's protocol is sensor-agnostic, but deployment requires a concrete
hardware reference. We specify a modular sensor pod that balances FOV
completeness, redundancy, time synchronization, and edge compute.

\subsection{Design Principles}

The platform is governed by four requirements: (1) $\ge 270^\circ$ horizontal
FOV across modalities; (2) at least two modalities covering every spatial
sector; (3) $<$1~ms inter-sensor time synchronization; and (4) all computation
(CT odometry, 3DGS map, compression) running locally on the edge GPU.

\subsection{Reference Platform}

The sensor pod is a self-contained, weatherproof (IP65) enclosure housing
all sensors, compute, power regulation, and communication.

\begin{table}[t]
  \centering
  \caption{Reference hardware platform.}
  \label{tab:hardware_spec}
  \footnotesize
  \setlength{\tabcolsep}{3pt}
  \begin{tabular}{@{}p{2.2cm}p{2.0cm}p{3.5cm}@{}}
    \toprule
    Component & Model/Spec & Notes \\
    \midrule
    Compute    & Jetson Orin AGX      & 2048-core Ampere, 32 GB, 275 TOPS \\
    Cameras $\times$4 & Alvium 1800 U-507c & 5 MP, global shutter, 30 Hz \\
    LiDARs $\times$2  & Livox Mid-360        & non-repetitive, 40$^\circ$ FOV \\
    IMU        & ICM-20948 + ADIS16470 & dual fused, 200 Hz \\
    Lens       & 8 mm, 82$^\circ$ HFOV  & C-mount, 30$^\circ$ overlap \\
    Wireless   & AX210 + X65             & Wi-Fi 6/6E + 5G~NR \\
    Time sync  & u-blox ZED-F9P          & GPS+PPS, cam trigger \\
    Storage    & 1 TB NVMe SSD           & submap buffer + logging \\
    Power      & 12--24 V DC, 45 W typ.\@ & PoE or 4S LiPo battery \\
    Enclosure  & CNC aluminum, IP65     & 220$\times$180$\times$140 mm, 2.8 kg \\
    \bottomrule
  \end{tabular}
\end{table}

\subsection{Sensor Layout}

Four cameras face outward at 90$^\circ$ intervals; with 82$^\circ$ HFOV each,
adjacent pairs overlap by 37$^\circ$. Two LiDARs are yawed $\pm 45^\circ$,
yielding 310$^\circ$ combined coverage. The body frame $B_i$ is centered at
the IMU.

\begin{figure}[t]
  \centering
  \resizebox{\columnwidth}{!}{% Hardware sensor layout — top-down view of the reference sensor pod
% 4 cameras at 0°, 90°, 180°, 270°; 2 LiDARs at ±45°
\begin{tikzpicture}[
  font=\fontsize{6}{7.2}\selectfont,
  >=Stealth,
]

% Pod outline
\draw[fill=gray!3, draw=gray!50, rounded corners] (-1.5,-1.5) rectangle (1.5,1.5);

% Body frame center
\draw[fill=red!60!black] (0,0) circle (0.08);
\node[font=\tiny, red!60!black] at (0.2,0.15) {$B_i$};

% Camera 0° (front) — right side
\draw[fill=blue!15, draw=blue!60!black, thick, rounded corners=1pt] (1.1,-0.3) rectangle (1.4,0.3);
\draw[blue!25, fill=blue!3] (1.25,0) -- (2.5,1.2) -- (2.5,-1.2) -- cycle;
\node[font=\tiny, blue!60!black] at (2.2,0.7) {Front};

% Camera 90° (right) — top side
\draw[fill=blue!15, draw=blue!60!black, thick, rounded corners=1pt] (-0.3,1.1) rectangle (0.3,1.4);
\draw[blue!25, fill=blue!3] (0,1.25) -- (1.2,2.5) -- (-1.2,2.5) -- cycle;
\node[font=\tiny, blue!60!black] at (0.7,2.2) {Right};

% Camera 180° (rear) — left side
\draw[fill=blue!15, draw=blue!60!black, thick, rounded corners=1pt] (-1.4,-0.3) rectangle (-1.1,0.3);
\draw[blue!25, fill=blue!3] (-1.25,0) -- (-2.5,1.2) -- (-2.5,-1.2) -- cycle;
\node[font=\tiny, blue!60!black] at (-2.2,0.7) {Rear};

% Camera 270° (left) — bottom side
\draw[fill=blue!15, draw=blue!60!black, thick, rounded corners=1pt] (-0.3,-1.4) rectangle (0.3,-1.1);
\draw[blue!25, fill=blue!3] (0,-1.25) -- (1.2,-2.5) -- (-1.2,-2.5) -- cycle;
\node[font=\tiny, blue!60!black] at (0.7,-2.2) {Left};

% LiDAR 1 at +45° (upper-right)
\draw[fill=green!15, draw=green!50!black, thick] (0.7,0) circle (0.22);
\node[font=\tiny, green!50!black] at (0.7,-0.35) {LiDAR};
\draw[green!20, fill=green!3] (0.7,0) -- (2.2,1.8) -- (0.1,2.5) -- cycle;

% LiDAR 2 at -45° (lower-left)
\draw[fill=green!15, draw=green!50!black, thick] (-0.7,0) circle (0.22);
\node[font=\tiny, green!50!black] at (-0.7,-0.35) {LiDAR};
\draw[green!20, fill=green!3] (-0.7,0) -- (-2.2,1.8) -- (-0.1,2.5) -- cycle;

\end{tikzpicture}
}
  \caption{Top-down sensor layout. Four cameras (blue cones) at 0$^\circ$/90$^\circ$/
  180$^\circ$/270$^\circ$; two LiDARs (green) at $\pm$45$^\circ$. $B_i$ at center.}
  \label{fig:layout}
\end{figure}

\subsection{Time Synchronization}

A u-blox ZED-F9P provides GPS UTC time (UART) and 1~Hz PPS ($<$20~ns jitter).
The PPS fans out to: (a) LiDAR sync inputs for point-cloud timestamping;
(b) Orin GPIO (\texttt{/dev/pps0}) for kernel-level timestamping via
\texttt{chrony} ($<10\,\mu$s accuracy); and (c) an STM32G0 trigger
generator producing 30~Hz camera pulses phase-locked to PPS. All sensor
timestamps are expressed in a common clock domain to within $<$1~ms.

\begin{figure}[t]
  \centering
  \resizebox{\columnwidth}{!}{% Time synchronization architecture diagram
% GPS → PPS → fan-out to LiDAR, Orin kernel, camera trigger generator
\begin{tikzpicture}[
  font=\fontsize{6}{7.2}\selectfont,
  >=Stealth,
  node distance=0.4cm and 0.6cm,
  box/.style={draw, rounded corners, minimum width=2.0cm, minimum height=0.45cm, align=center, inner sep=2pt},
  sync/.style={->, very thick},
]

% GPS module
\node[box, fill=yellow!12] (gps) {u-blox ZED-F9P (GPS+PPS)};

% PPS fan-out
\node[box, fill=orange!12, below=0.8cm of gps] (pps) {PPS Fan-Out (1 Hz)};

% LiDAR
\node[box, fill=green!12, right=1.0cm of pps, yshift=0.5cm] (lidar) {LiDARs $\times$2 Mid-360};

% Orin GPIO (middle row)
\node[box, fill=blue!10, right=1.0cm of pps, yshift=-0.5cm] (orin) {Orin GPIO PPS};

% Camera trigger gen
\node[box, fill=red!10, right=0.8cm of pps] (trig) {STM32G0 Trig.\@ Gen.\@ (30 Hz)};

% Cameras
\node[box, fill=gray!10, right=1.0cm of trig, minimum height=0.9cm] (cams) {Cameras $\times$4 Alvium};

% Arrows
\draw[sync] (gps.south) -- (pps.north);
\draw[sync, green!50!black] (pps.east) -- (lidar.west);
\draw[sync, blue!60!black] (pps.east) -- (orin.west);
\draw[sync, red!60!black] (pps.east) -- (trig.west);
\draw[sync] (trig.east) -- (cams.west);

% Chrony
\node[box, fill=purple!8, below=0.8cm of orin] (chrony) {chrony (PTP)};
\draw[sync, purple!60!black, dashed] (orin.south) -- (chrony.north);

\end{tikzpicture}
}
  \caption{Time synchronization architecture. GPS+PPS $\to$ fan-out to
  LiDAR/Orin/camera trigger generator. All timestamps in PPS-aligned domain.}
  \label{fig:time_sync}
\end{figure}

\subsection{Power and Heat}

Steady-state consumption is 45~W (Orin: 15~W, LiDARs: 12~W, cameras: 5~W,
wireless: 3~W, GNSS+MCU: 3~W, misc: 7~W). Under 35$^\circ$C ambient, Orin
GPU junction temperature stabilizes at 72$^\circ$C (below the 85$^\circ$C
throttling threshold) with forced-air cooling at 3000~RPM.

For UAV deployment, a 1.1~kg, 22~W \textbf{UAV-lite} variant uses 1 camera
+ 1 LiDAR + Orin NX (8~GB). Full platform details are encoded in the node
descriptor (Section~\ref{sec:descriptor}) to inform cluster-head election
and Tier-0 inventory.
